% Статья 2 из 2 (разделение по топологии). "Точная вероятность нарушения SLA
% в многоканальных системах пакетного обслуживания с общей очередью:
% нетерпеливые заявки".
%
% Статья 1 (одноканальная топология, базовая конструкция + H2/Ek, EPIC-035/036/
% 042/067/068) -- в соседнем каталоге ../batch-service-sla-single/latex/main.tex.
%
% Личный материал (не часть пакета most-queue, docs/статьи/ в .gitignore).
% Стиль оформления -- по образцу реальной опубликованной статьи автора в
% журнале "Труды Военно-космической академии имени А.Ф. Можайского".
%
% Сборка: latexmk -xelatex main.tex   (требуется XeLaTeX + biber)
%
% encoding: UTF-8
\documentclass[12pt,a4paper]{article}

\usepackage{geometry}
\geometry{top=2.5cm,bottom=2.5cm,left=3cm,right=2cm}

\usepackage{fontspec}
\setmainfont{PT Serif}
\newfontfamily\cyrillicfonttt{DejaVu Sans Mono}

\usepackage{polyglossia}
\setmainlanguage{russian}
\setotherlanguage[variant=american]{english}

\usepackage{amsmath,amssymb}
\usepackage{graphicx}
\usepackage{siunitx}
\sisetup{output-decimal-marker = {,}, group-separator = {\,}, retain-unity-mantissa = false}

\usepackage{array}
\usepackage{multirow}
\usepackage{caption}
\usepackage{capt-of}
\usepackage{seqsplit}
\usepackage{xcolor}
\usepackage{tikz}
\usetikzlibrary{arrows.meta, positioning, shapes.geometric}
\captionsetup[figure]{labelsep=period, justification=centering, font=small, position=bottom}
\captionsetup[table]{labelsep=period, font={small,bf}, justification=raggedright, position=top}

\usepackage{indentfirst}
\setlength{\parindent}{1.25cm}
\linespread{1.12}

\usepackage[backend=biber, style=gost-numeric, sorting=none]{biblatex}
\addbibresource{../../bibliography/library.bib}

% Заголовки разделов: без нумерации, жирным, без капслока -- по образцу
% реальной статьи автора в "Трудах ВКА" (не ГОСТовский численный рубрикатор).
\newcommand{\articlesection}[1]{\par\vspace{14pt}{\bfseries #1}\par\vspace{6pt}\nopagebreak}

\newcommand{\keywordslabel}{\textbf{Ключевые слова: }}

\begin{document}
\sloppy

\noindent УДК 519.872

\bigskip
\noindent Р.С.~Хабаров

\noindent{\small [степень, звание, место работы --- заполнить]}

\bigskip
\begin{center}
\begin{minipage}{0.9\linewidth}
\centering
{\large\bfseries ТОЧНАЯ ВЕРОЯТНОСТЬ НАРУШЕНИЯ SLA В МНОГОКАНАЛЬНЫХ СИСТЕМАХ ПАКЕТНОГО ОБСЛУЖИВАНИЯ С ОБЩЕЙ ОЧЕРЕДЬЮ: НЕТЕРПЕЛИВЫЕ ЗАЯВКИ}
\end{minipage}
\end{center}

\bigskip
\noindent\small
Рассмотрена система из $n$ идентичных каналов пакетного (batch-service) обслуживания с общей очередью и пороговым запуском пакета --- модель пула из нескольких реплик GPU-инференса за общим маршрутизатором запросов. Предложен точный метод вычисления вероятности нарушения SLA $P(W>D)$ для непрерывного времени, произвольных $n\ge1$ и $1\le a\le b$: при батч-размер-независимом обслуживании агрегированная ставка освобождения любого из $n$ занятых каналов равна ровно $n\mu$ независимо от состава их занятости, что сводит ожидание помеченной заявки к уже решённой одноканальной конструкции с масштабированной ставкой. Метод распространён на марковское нетерпение заявок (та же агрегированная ставка переносит одноканальную конструкцию без изменений) и на батч-размер-\textit{зависимое} обслуживание, для которого это упрощение перестаёт выполняться буквально: состав занятости каналов отслеживается внутри поглощающей цепи явно, но лишь пока заняты все $n$ каналов, что удерживает пространство состояний заметно меньше первоначально ожидавшегося. Все результаты точно регрессируют к одноканальному случаю, к отсутствию нетерпения и к постоянному обслуживанию, и подтверждены независимым имитационным моделированием. Численно показано: добавление одной GPU-реплики в предельно нагруженный пул снижает среднее время ожидания более чем в 17~раз и вероятность нарушения дедлайна --- более чем на порядок; нетерпение заявок и дополнительная ёмкость пула --- взаимозаменяемые средства снижения ожидания, с убывающей отдачей от нетерпения по мере роста базовой ёмкости; а приближение зависящего от размера обслуживания одной постоянной ставкой систематически завышает и среднее время ожидания, и, заметно сильнее, вероятность нарушения дедлайна в хвосте.

\medskip
\noindent\keywordslabel теория массового обслуживания, пакетное обслуживание, многоканальные системы, общая очередь, SLA, вероятность нарушения дедлайна, нетерпеливые заявки, обслуживание, зависящее от размера пакета, инференс больших языковых моделей, динамический батчинг, пул GPU-реплик.
\normalsize

\bigskip
\begin{center}
\begin{minipage}{0.9\linewidth}
\centering
{\bfseries EXACT SLA-VIOLATION PROBABILITY FOR MULTI-SERVER THRESHOLD BATCH-SERVICE QUEUES WITH A SHARED QUEUE: CUSTOMER ABANDONMENT}
\end{minipage}
\end{center}

\begin{center}
\small R.S.~Khabarov
\end{center}

\begin{otherlanguage}{english}
\small
We consider a pool of $n$ identical batch-service servers sharing one queue with a threshold batch-formation rule --- the natural model for a pool of GPU-inference replicas behind a shared request router. We derive an exact deadline-violation probability $P(W>D)$ for continuous time, any $n\ge1$ and $1\le a\le b$: under batch-size-independent service, the aggregate rate at which any of the $n$ busy servers next completes a batch is exactly $n\mu$ regardless of occupancy composition, which reduces the tagged customer's wait to the already-solved single-server construction with a rescaled rate. The method is extended to Markovian customer abandonment (the same aggregate rate carries the single-server construction over unchanged) and to batch-size-DEPENDENT service, for which this simplification no longer holds literally: occupancy composition is tracked explicitly inside the tagged customer's absorbing chain, but only while all $n$ servers stay busy, which keeps the actual state space considerably smaller than initially expected. All results regress exactly to the single-server case, to no abandonment, and to constant service, and are validated against independent discrete-event simulation. Numerically: adding one GPU replica to a saturated pool cuts the mean wait by more than 17-fold and the deadline-violation probability by more than an order of magnitude; customer abandonment and extra capacity are substitutes for reducing wait, with diminishing returns from abandonment as baseline capacity grows; and approximating size-dependent service with a single constant rate systematically overstates both the mean wait and, considerably more so, the tail deadline-violation probability.

\medskip
\noindent\textit{Key words:} queueing theory, batch service, multi-server systems, shared queue, SLA, deadline-violation probability, customer abandonment, batch-size-dependent service, large language model inference, dynamic batching, GPU replica pool.
\end{otherlanguage}
\normalsize

\renewcommand{\figurename}{Рисунок}
\renewcommand{\tablename}{Таблица}

\articlesection{Введение}

Ускорение инференса нейронных сетей на графических процессорах (GPU) существенно опирается на динамический батчинг --- объединение нескольких независимо поступивших запросов в один пакет (batch), обрабатываемый каналом одновременно; такое объединение повышает пропускную способность канала ценой дополнительной задержки ожидания формирования и обслуживания пакета \cite{inoue2021}. Операторы промышленных систем инференса, однако, почти никогда не ограничиваются одной GPU: типичное развёртывание --- пул из нескольких идентичных реплик модели за общим маршрутизатором запросов, распределяющим нагрузку по общей очереди. Для таких пулов требование к качеству обслуживания формулируется в терминах вероятности, а не среднего: $P(\mathrm{TTFT}>D)<\varepsilon$ для времени до первого токена и заданного малого $\varepsilon$ \cite{quartz2026,jitserve2026}, что делает точное (а не приближённое по среднему) распределение времени ожидания практически значимым и для многоканальной топологии.

Точное стационарное распределение длины очереди для систем пакетного обслуживания с зависящим от размера пакета временем обслуживания детально проработано, начиная с работы \cite{banerjee2012} и в серии последующих работ \cite{pradhan2016,banerjee2015,gupta2018,pradhan2024}, но для одноканальной топологии; точное время ожидания помеченной заявки для общего (не обязательно зависящего от размера) обслуживания получено методом корней характеристического уравнения \cite{chaudhry2023} --- также только для одного канала. Автор этих строк ранее точно решил одноканальный случай с зависящим от размера пакета фазовым обслуживанием и его сочетание с нетерпением заявок \cite{khabarovsingle2026}; настоящая работа распространяет точный результат на топологию с несколькими параллельными каналами и общей очередью.

Несколько параллельных каналов батч-обслуживания изучались и ранее: недавняя работа \cite{bruneel2026} даёт точный анализ двух параллельных очередей с геометрическим батч-обслуживанием, но в дискретном времени и без общей очереди --- каждая из двух очередей обслуживается своим каналом независимо, без разделения нагрузки между каналами. Это принципиальное отличие от рассматриваемой здесь топологии: общая очередь означает, что любой из $n$ каналов может обслужить следующую по очереди заявку, и именно это разделение нагрузки определяет выигрыш от масштабирования числа каналов, количественно исследованный ниже. Нетерпение заявок в сочетании с несколькими каналами пакетного обслуживания и общей очередью ранее точно, по нашим сведениям, не рассматривалось.

Дальнейшая структура статьи следующая. Сначала приводится постановка задачи. Далее описан метод: конструкция для $n$ идентичных каналов с общей очередью и её расширение на нетерпение заявок. После этого представлены численные эксперименты, количественно оценивающие эффект масштабирования числа каналов и его взаимодействие с нетерпением заявок.

\articlesection{Постановка задачи}

Рассматривается система $M/M^{[a,b]}/n$: пуассоновский входной поток интенсивности $\lambda$; $n\ge1$ идентичных каналов обслуживают пакеты заявок из одной общей очереди; каждый освободившийся канал начинает обслуживание пакета, как только накопится не менее $a$ ожидающих заявок, и забирает в пакет $\min(b,j)$ заявок, где $j$ --- число ожидающих в момент запуска ($1\le a\le b$ --- общий порог запуска). Время обслуживания пакета размера $s$ --- экспоненциальное с темпом $\mu(s)$, одинаковым для всех $n$ каналов; базовый случай --- постоянная ставка $\mu(s)\equiv\mu$, не зависящая от размера пакета.

Базовая цель --- точное $P(W>D)$ для $D\ge0$, где $W$ --- время ожидания в очереди случайно выбранной (по PASTA) заявки, при любом $n\ge1$ и любом $1\le a\le b$, при батч-размер-независимом обслуживании. Эта конструкция расширяется в двух направлениях. Во-первых, на нетерпение заявок: каждая из $j$ ожидающих заявок независимо покидает очередь с постоянной ставкой $\gamma$ (марковское/беспамятное терпение, конвенция модели Erlang-A \cite{garnett2002}), если её обслуживание не началось; цель --- точная вероятность того, что помеченная заявка покинет очередь до начала своего пакета, и точное условное (на <<дождалась>>) распределение $W$. Во-вторых, на батч-размер-\textit{зависимое} обслуживание $\mu(s)$, произвольную функцию размера пакета; цель --- то же точное $P(W>D)$, но уже без батч-размер-независимого упрощения, на котором держится базовая конструкция. Оба расширения получены раздельно (их сочетание --- нетерпение при batch-size-зависимом обслуживании --- остаётся резервом, см. раздел <<Ограничения>>).

\articlesection{Метод}

\noindent\textbf{Базовая конструкция, $n$ идентичных каналов.} Состояние системы --- вектор занятости $(m_1,\dots,m_b)$ (число каналов, занятых пакетом размера $s$, $s=1,\dots,b$; $n-\sum m_s$ --- число свободных каналов), скрещенный с числом ожидающих заявок $j$. Пакет при хотя бы одном свободном канале запускается немедленно по достижении порога $a$ (той же конвенцией, что и в одноканальном случае \cite{khabarovsingle2026}) --- отсюда встроенный инвариант: свободный канал есть $\Rightarrow j<a$ (иначе он уже забрал бы заявки). Стационарное распределение находится прямым решением $\pi Q=0$ для разреженного генератора $Q$ этой цепи Маркова.

Ключевое упрощение, не предполагавшееся заранее, --- для батч-размер-независимого $\mu$ агрегированная ставка <<следующее освобождение>> среди всех занятых каналов равна ровно (число занятых каналов)$\cdot\mu$ независимо от состава их занятости (по беспамятности экспоненциального обслуживания и немедленному повторному запуску освободившегося канала, если накопленных заявок достаточно). В частности, когда заняты все $n$ каналов, эта ставка равна ровно $n\mu$. Время ожидания помеченной заявки при всех занятых каналах поэтому не требует отслеживания вектора занятости вовсе: число заявок впереди помеченной убывает на $b$ при каждом событии освобождения со ставкой $n\mu$ --- в точности одноканальная конструкция \cite{khabarovsingle2026} с единой ставкой $n\mu$ на всём протяжении вместо исходной ставки $\mu$ (и без различия <<первый пакет vs последующие>>, поскольку здесь нет понятия <<застигнутый пакет>>: любой из $n$ одинаковых каналов может освободиться следующим). Если хотя бы один канал свободен, $W=0$ (помеченная заявка сразу попадает в пакет на свободном канале) либо время ожидания сводится к ожиданию накопления недостающих заявок по пуассоновскому потоку -- та же конструкция, что и в одноканальном случае.

\noindent\textbf{Расширение: нетерпение заявок.} Каждая из $j$ ожидающих заявок покидает очередь со ставкой $\gamma$. Благодаря тому же упрощению агрегированной ставки, одноканальная конструкция нетерпения \cite{khabarovsingle2026} --- явный конечный подгенератор $T$ поглощающей цепи Маркова на состоянии $(R,K)$ ($R$ --- число выживших впереди помеченной заявки, $K$ --- новые приходы позади неё), с конкурирующими поглощающими типами <<обслужена>> (ставка освобождения канала) и <<ушла по терпению>> (ставка $\gamma$, пропорциональная числу выживших впереди), --- переносится без каких-либо изменений в самой структуре состояния: единственное отличие --- подстановка ставки обслуживания $n\mu$ вместо $\mu$ там, где конструкция для одного канала использует ставку освобождения занятого канала. Формулы вероятности ухода и условных моментов времени ожидания --- те же, что и в одноканальном случае:
\begin{equation}
P(\text{уход}) = \gamma\,\alpha\,(-T)^{-1}\mathbf{1}, \qquad
\mathrm{E}[W^k;\text{обслужена}] = k!\,\alpha\,(-T)^{-(k+1)}\,r,
\label{eq:abandonment}
\end{equation}
нормированные на $P(\text{обслужена})=1-P(\text{уход})$ для получения условных моментов; $r$ --- вектор ставок обслуживания, здесь равный $n\mu$ в каждой компоненте, отвечающей состоянию, где формируется пакет, включающий помеченную заявку.

Регрессионные тождества --- по $n=1$ (точное совпадение с одноканальной конструкцией \cite{khabarovsingle2026}, включая случай $\gamma>0$) и по $\gamma=0$ (точное совпадение с конструкцией без нетерпения) --- выполняются точно. Независимая проверка имитационным моделированием (пул из $n$ каналов с общей очередью, каждая ожидающая заявка независимо покидает очередь со ставкой $\gamma$; время ожидания измерено как момент начала обслуживания минус момент прихода, а не момент завершения пакета) подтверждает точность конструкции на нескольких комбинациях $(a,b,n)$ и значений $\gamma$, в том числе по вероятности ухода и условному среднему времени ожидания.

\noindent\textbf{Расширение: батч-размер-зависимое обслуживание.} Если темп $\mu(s)$ зависит от размера пакета $s$, ключевое упрощение базовой конструкции перестаёт выполняться буквально: агрегированная ставка <<следующее освобождение>> среди занятых каналов равна $\sum_s n_s\,\mu(s)$, где $n_s$ --- число каналов, занятых пакетом размера $s$, --- и генуинно зависит от состава занятости $(n_1,\dots,n_b)$, а не только от их числа. Поглощающая цепь помеченной заявки поэтому несёт состав занятости как явную часть состояния: $(\text{состав},R,K)$, где $R$ и $K$ --- те же величины, что и в расширении на нетерпение. Каждый занятый класс размера $s$ (с $n_s>0$) вносит собственную конкурирующую ставку $n_s\,\mu(s)$; при срабатывании --- если накопленного $R+1+K$ всё ещё достаточно для немедленного повторного запуска ($\ge a$) и это не включает помеченную заявку --- состав предсказуемо обновляется ($n_s\mathbin{-{=}}1$, $n_{\min(b,R+1+K)}\mathbin{+{=}}1$, поскольку освободившийся канал немедленно забирает новый пакет), и цепь продолжается с уменьшенным $R$.

Ключевая находка, не предполагавшаяся заранее и заметно сужающая ожидавшийся взрыв размерности состояния, --- состав занятости нужно отслеживать \textit{только} пока заняты все $n$ каналов. В момент, когда освободившийся канал не может сразу забрать новый пакет (накопленного $R+1+K$ всё ещё меньше $a$), он простаивает --- и с этого момента дальнейшие завершения остальных ещё занятых каналов уже не способны изменить содержимое очереди (канал, завершающий обслуживание при неизменном числе ожидающих, лишь создаёт ещё один простаивающий канал, а не диспетчеризацию). Поэтому с этого момента действует уже существующая, не зависящая от состава занятости гонка ожидания накопления по пуассоновскому потоку --- та же конструкция, что и в расширении на нетерпение (при $\gamma=0$). Это ограничивает отслеживаемое пространство составов занятости только векторами с суммой ровно $n$ (а не $\le n$, как для полного стационарного распределения) --- заметно меньшее множество, не зависящее от усечения очереди.

На рисунке~\ref{fig:occupancy-diagram} приведён пример эволюции состава занятости для $n=2$ каналов, $b=2$ (темпы $\mu(1)$, $\mu(2)$), показывающий только переходы, сохраняющие оба канала занятыми (случай <<ещё достаточно впереди>>; выходы в <<обслужена>> и <<простой>> зависят от $R$, а не от состава, и здесь не изображены --- та же логика, что и для поглощения в расширении на нетерпение).

\begin{center}
\resizebox{0.62\linewidth}{!}{%
\begin{tikzpicture}[
  every node/.style={font=\footnotesize},
  state/.style={ellipse, draw=black, thick, minimum width=16mm, minimum height=9mm, inner sep=0pt, align=center},
  occ/.style={state, fill=blue!12},
  ph/.style={-{Stealth[length=2.2mm]}, thick, black},
]

\node[occ] (s20) at (0,2.2) {$(2,0)$};
\node[occ] (s11) at (5.0,2.2) {$(1,1)$};
\node[occ] (s02) at (10.0,2.2) {$(0,2)$};

\draw[ph] (s20) -- node[above]{$2\mu(1)$} (s11);
\draw[ph] (s11) -- node[above]{$\mu(1)$} (s02);
\draw[ph] (s11) to[out=60,in=120,looseness=8] node[above]{$\mu(2)$} (s11);
\draw[ph] (s02) to[out=60,in=120,looseness=8] node[above]{$2\mu(2)$} (s02);

\end{tikzpicture}}
\captionof{figure}{Эволюция состава занятости $(n_1,n_2)$ для $n=2$ каналов, $b=2$ (только переходы, сохраняющие оба канала занятыми); узел --- число каналов, занятых пакетом размера 1 и размера 2 соответственно}
\label{fig:occupancy-diagram}
\end{center}

Регрессионное тождество --- по постоянному $\mu(s)\equiv\mu$ (точное, до полной точности float64, совпадение с базовой конструкцией) --- выполняется точно. Независимая проверка имитационным моделированием на генуинно батч-размер-зависимом $\mu(s)=1/(0{,}3+0{,}25s)$ (таблица~\ref{tab:ryzhikov-dependent-mu}) подтверждает точность конструкции по среднему времени ожидания и хвосту $P(W>t)$.

\begin{center}
\captionof{table}{Расчёт и имитация: batch-size-зависимое обслуживание ($2{,}5\times10^6$ заявок на прогон)}
\label{tab:ryzhikov-dependent-mu}
\begin{tabular}{|c|c|c|c|c|c|c|}
\hline
\multirow{2}{*}{$(a,b,n,\lambda)$} & \multicolumn{3}{c|}{$\mathrm{E}[W]$} & \multicolumn{3}{c|}{$P(W>1)$} \\
\cline{2-7}
 & теор. & ИМ & отн. ошибка & теор. & ИМ & отн. ошибка \\
\hline
$(2,4,2,1{,}0)$ & 0,5072 & 0,5063 & 0,18\,\% & 0,1861 & 0,1860 & 0,04\,\% \\
$(2,3,2,1{,}3)$ & 0,4013 & 0,4006 & 0,18\,\% & 0,1416 & 0,1414 & 0,16\,\% \\
\hline
\end{tabular}
\end{center}

Расхождение не превышает 0,18\,\% по среднему времени ожидания и 0,16\,\% по хвосту --- в пределах статистической погрешности, ожидаемой при таком объёме испытаний.

\articlesection{Численные эксперименты}

\noindent\textbf{Эффект числа каналов.} На параметрах batch-size-зависимости среднего \cite{inoue2021} ($\alpha=0{,}1438$, $\tau_0=1{,}8874$), с батч-размер-независимым экспоненциальным обслуживанием со средним, равным среднему при размере пакета 4 ($\mu=1/(4\alpha+\tau_0)\approx0{,}4061$ --- максимальная пропускная способность одного канала при $b_{\max}=4$ составляет $4\mu\approx1{,}624$), $a=2$, $b_{\max}=4$, общая интенсивность входного потока $\lambda=1{,}5$ (97,5\,\% от максимальной пропускной способности одного канала --- система с $n=1$ стабильна, но предельно нагружена), таблица~\ref{tab:multiserver} показывает эффект добавления идентичных каналов в общий пул с общей очередью.

\begin{center}
\captionof{table}{Эффект числа каналов $n$ при фиксированной суммарной нагрузке ($\lambda=1{,}5$, $\gamma=0$)}
\label{tab:multiserver}
\begin{tabular}{|c|c|c|c|c|}
\hline
$n$ & $\mathrm{E}[W]$ & $P(W>5)$ & $P(W>10)$ & Загрузка канала \\
\hline
1 & 19,754 & 0,7633 & \num{5.97e-1} & 0,975 \\
2 & 1,144  & \num{2.95e-2} & \num{1.45e-3} & 0,739 \\
3 & 0,509  & \num{1.47e-3} & \num{5.63e-6} & 0,574 \\
4 & 0,370  & \num{3.15e-4} & \num{1.71e-7} & 0,454 \\
\hline
\end{tabular}
\end{center}

Переход от одного предельно нагруженного канала к двум ($n=1\to2$) снижает $\mathrm{E}[W]$ более чем в 17~раз и $P(W>5)$ --- более чем в 25~раз, хотя суммарная входная нагрузка не изменилась: с одним каналом система работает на грани устойчивости ($\rho=97{,}5\,\%$), и каждая дополнительная GPU-реплика пула не просто делит нагрузку поровну, а снимает именно ту часть перегрузки, которая отвечает за формирование длинного хвоста ожидания. При дальнейшем увеличении $n$ отдача убывает (переход $n=2\to3$ даёт более чем 20-кратное улучшение $P(W>5)$, а переход $n=3\to4$ --- <<лишь>> примерно 4,7-кратное) --- типичная для систем массового обслуживания картина убывающей отдачи от масштабирования числа каналов, здесь впервые точно, а не оценочно, количественно определённая для системы пакетного обслуживания с общей очередью. Практический вывод для оператора GPU-пула: при близкой к насыщению нагрузке одна дополнительная реплика значительно ценнее, чем может показаться по простому <<делению нагрузки на $n$>>.

\noindent\textbf{Взаимодействие нетерпения заявок с числом каналов.} На тех же параметрах (таблица~\ref{tab:abandonment-multiserver}) показан эффект ставки терпения $\gamma$ при трёх разных $n$ --- вероятность отмены запроса до начала его пакета $P(\text{уход})$, условное (на <<дождался>>) среднее время ожидания $\mathrm{E}[W\mid\text{обслужена}]$ и истинная вероятность выполнения SLA с дедлайном $D=10$ --- $P(\text{обслужена и }W\le D)$.

\begin{center}
\captionof{table}{Нетерпение заявок при разном числе каналов ($D=10$)}
\label{tab:abandonment-multiserver}
\begin{tabular}{|c|c|c|c|c|}
\hline
$n$ & $\gamma$ & $P(\text{уход})$ & $\mathrm{E}[W\mid\text{обслужена}]$ & $P(\text{обслужена и }W\le D)$ \\
\hline
1 & 0    & 0      & 19,93 & 0,4022 \\
1 & 0,1  & 0,2496 & 2,51  & 0,7388 \\
1 & 0,3  & 0,4011 & 1,27  & 0,5984 \\
\hline
2 & 0    & 0      & 1,14 & 0,9985 \\
2 & 0,1  & 0,0864 & 0,84 & 0,9135 \\
2 & 0,3  & 0,1914 & 0,59 & 0,8086 \\
\hline
3 & 0    & 0      & 0,509 & 0,99999 \\
3 & 0,1  & 0,0459 & 0,446 & 0,95410 \\
3 & 0,3  & 0,1185 & 0,365 & 0,88148 \\
\hline
\end{tabular}
\end{center}

Таблица~\ref{tab:abandonment-multiserver} вскрывает качественно различную картину в зависимости от $n$. При $n=1$ (предельно нагруженный одиночный канал, $\rho=97{,}5\,\%$) нетерпение заявок --- не только статистический артефакт, но и ощутимый практический рычаг: оно само по себе действует как предохранительный клапан, отсекающий избыточную нагрузку и поднимающий истинную вероятность выполнения SLA с 40,2\,\% ($\gamma=0$, система фактически не справляется с нагрузкой на горизонте $D=10$) до 73,9\,\% уже при скромной ставке $\gamma=0{,}1$ --- рост, целиком проистекающий из статистического отбрасывания части нагрузки, а не из улучшения самой системы; при дальнейшем росте $\gamma$ выигрыш не монотонен (при $\gamma=0{,}3$ показатель снижается до 59,8\,\%), поскольку увеличивающаяся доля уходов в конце концов начинает учитываться как провал SLA наравне с нарушением дедлайна. При $n=3$ (загрузка канала уже умеренная, 57,4\,\%) та же ставка терпения $\gamma=0{,}3$ отсекает заметно меньшую долю нагрузки (11,85\,\% против 40,1\,\% при $n=1$) и из уже близкой к единице вероятности SLA-успеха без нетерпения (99,999\,\%) снижает её до 88,15\,\% --- то есть при низкой базовой загрузке нетерпение заявок способно только ухудшить наблюдаемый SLA-показатель, отбрасывая заявки, которые и так были бы обслужены в срок.

Это --- количественное проявление взаимозаменяемости двух средств снижения времени ожидания, упомянутой во введении: нетерпение заявок компенсирует недостаток ёмкости пула, но ценой прямой потери части нагрузки, тогда как дополнительная ёмкость устраняет саму причину долгого ожидания без этой потери. Для оператора, управляющего пулом GPU-реплик, вывод прямой: наблюдаемое улучшение SLA-метрик при увеличении ставки нетерпения (или при агрессивных клиентских таймаутах, эквивалентных росту $\gamma$) не свидетельствует об улучшении системы и не должно восприниматься как альтернатива добавлению ёмкости при хронической перегрузке пула.

\noindent\textbf{Цена приближения постоянным обслуживанием.} Таблица~\ref{tab:multiserver} выше вычислена при единственной ставке $\mu=1/(4\alpha+\tau_0)$ --- приближении, соответствующем размеру пакета $b_{\max}=4$, то есть предположению, что каждый пакет полноразмерный. На тех же параметрах, но с истинным батч-размер-зависимым $\mu(s)=1/(\alpha s+\tau_0)$ (та же подгонка по среднему, что и всюду в статье), таблица~\ref{tab:dependent-mu} сравнивает точный результат с этим приближением при $\gamma=0$.

\begin{center}
\captionof{table}{Точное батч-размер-зависимое обслуживание против приближения постоянной ставкой $\mu(b_{\max})$}
\label{tab:dependent-mu}
\begin{tabular}{|c|c|c|c|c|}
\hline
$n$ & Точное $\mathrm{E}[W]$ & Приближение & Точное $P(W>10)$ & Приближение \\
\hline
2 & 0,964 & 1,144 & \num{8.04e-4} & \num{1.45e-3} \\
3 & 0,448 & 0,509 & \num{1.45e-6} & \num{5.63e-6} \\
\hline
\end{tabular}
\end{center}

Приближение постоянной ставкой систематически \textit{завышает} и среднее время ожидания (на 18,7\,\% при $n=2$, на 13,6\,\% при $n=3$), и, заметно сильнее, вероятность нарушения дедлайна в хвосте (на 80\,\% при $n=2$, почти в 4~раза при $n=3$) --- направление ошибки объясняется тем, что реальные пакеты зачастую меньше $b_{\max}$ (особенно вблизи порога $a$), а при убывающей по размеру ставке $\mu(s)=1/(\alpha s+\tau_0)$ меньшие пакеты обслуживаются быстрее: подстановка темпа самого медленного (полноразмерного) пакета вместо истинного смешанного темпа делает пул искусственно более перегруженным, чем он есть на самом деле. Эффект усиливается с глубиной хвоста --- та же качественная картина, что и для однополюсного приближения хвоста в \cite{khabarovsingle2026}: грубое приближение одной ставкой терпимо для среднего, но ненадёжно именно там, где нужна точность SLA-ориентированного управления.

\articlesection{Ограничения}

Батч-размер-зависимое обслуживание и нетерпение заявок реализованы раздельно; их сочетание остаётся резервом --- поглощающая цепь состава занятости (раздел <<Метод>>) не несёт переходов по уходу заявок впереди, перенос которых потребовал бы их добавления к уже и так расширенному состоянию. Фазовое (Erlang/$H_2$) обслуживание при $n>1$ также оставлено резервом: оно добавило бы фазовое измерение к каждому занятому каналу поверх уже отслеживаемого состава занятости. Нетерпение заявок для $H_2$-обслуживания даже в одноканальном случае --- отдельный резерв, рассмотренный в \cite{khabarovsingle2026}, не затрагиваемый здесь.

Численные примеры (таблицы~\ref{tab:multiserver}--\ref{tab:dependent-mu}) --- иллюстративные, на одном наборе параметров ($a=2$, $b_{\max}=4$, фиксированная суммарная нагрузка) каждый; систематическое исследование по сетке параметров (аналогичное исследованию порога $a$ и $b_{\max}$ в \cite{khabarovsingle2026}) для многоканальной топологии не проводилось. В частности, наблюдение о немонотонности эффекта нетерпения по $n$ (ощутимая польза при высокой загрузке, чистый вред при низкой) --- качественный вывод на конкретных параметрах, а не доказанное в общем виде свойство; то же относится к направлению и величине ошибки приближения постоянным обслуживанием (таблица~\ref{tab:dependent-mu}), проверенного только для одной убывающей по размеру функции $\mu(s)$.

\articlesection{Заключение}

Получена и верифицирована точная вероятность нарушения SLA для системы пакетного обслуживания из $n$ идентичных каналов с общей очередью и общим порогом запуска пакета $1\le a\le b$, при батч-размер-независимом времени обслуживания, --- результат, основанный на точном показе того, что агрегированная ставка освобождения любого из занятых каналов не зависит от состава их занятости, что сводит многоканальную конструкцию к уже решённой одноканальной \cite{khabarovsingle2026} с масштабированной ставкой. Метод распространён на марковское нетерпение заявок, для которого та же агрегированная ставка позволяет перенести одноканальную конструкцию нетерпения без изменения её структуры, и на батч-размер-зависимое обслуживание, для которого это упрощение перестаёт выполняться буквально и состав занятости каналов приходится отслеживать явно внутри поглощающей цепи помеченной заявки --- но, как показано, только пока все каналы заняты, что удерживает фактическое пространство состояний заметно меньше того, что предполагал первоначальный анализ риска.

Численно показано, что добавление одной GPU-реплики в предельно нагруженный пул снижает среднее время ожидания более чем в 17~раз, а вероятность нарушения дедлайна --- более чем на порядок, с типичной для масштабирования числа каналов убывающей отдачей; что нетерпение заявок и дополнительная ёмкость пула действуют как в существенной мере взаимозаменяемые средства: та же ставка терпения резко повышает истинную вероятность выполнения SLA при высокой базовой загрузке пула, но при умеренной загрузке лишь отбрасывает часть нагрузки, которая и так была бы обслужена в срок, ухудшая наблюдаемый показатель; и что приближение батч-размер-зависимого обслуживания одной постоянной ставкой, откалиброванной по самому крупному размеру пакета, систематически завышает среднее время ожидания на 14--19\,\% и вероятность нарушения дедлайна в хвосте --- в разы. Эти результаты прямо предостерегают от использования клиентских таймаутов как замены достаточной ёмкости инференс-пула и от грубого приближения обслуживания одной ставкой там, где нужна точность SLA-ориентированного управления. Открытая программная реализация всех описанных методов доступна в библиотеке most-queue \cite{mostqueue2026}.

\articlesection{Список литературы}
\printbibliography[heading=none]

\end{document}
